\section{Proper equations and admissible approximation}
\label{sec:hessian-proper-approximation}
\index{Hessian equation!proper exponential equation}
\index{Admissible potential!smooth approximation}

For a smooth potential write
$H_m(w)=(\omega+\ddc w)^m\wedge\omega^{n-m}$.
The positive exponential coefficient controls additive constants by the
maximum principle. The resulting smooth equations will also construct
admissible approximations of continuous potentials, before any continuous
Hessian product is defined.

\Needspace{11\baselineskip}
\begin{lemma}\label{hessian:smooth-proper-existence}
For every $a>0$ and smooth real $H$ on $X$, there is a unique smooth
strictly admissible solution of
\[
 H_m(w)=e^{aw+H}\omega^n.
\]
It satisfies
\[
 -\frac{\max_X H}{a}\leq w\leq-\frac{\min_X H}{a}.
\]
\end{lemma}
\begin{proof}
At a maximum of $w$, the admissible form $\omega+\ddc w$ is at most
$\omega$, hence $e^{aw+H}\leq1$ there. At a minimum it is at least
$\omega$, and the reverse inequality holds. These give the bounds.
For two solutions, a positive maximum of their difference would imply
$H_m(w_1)\leq H_m(w_2)$ there, contradicting the exponential factor.
This proves uniqueness.

First take $a=1$ and use the path
$H_m(w_t)=e^{w_t+tH}\omega^n$, starting with $w_0=0$.
The trace estimate of \cref{hessian:prescribed-potential-trace} applies with prescribed potential $v=0$,
fixed smooth comparison $0$, and $h(x,v)=tH(x)$; its constants are
uniform for $t\in[0,1]$. \Cref{hessian:hessian-trace-to-gradient} then gives a
uniform gradient and complex-Hessian bound. \Cref{hessian:complex-hessian-compactness} gives uniform $W^{2,p}$ and hence $C^{1,\alpha}$ bounds, for any
fixed $\alpha<1$. Consequently
$r_t=e^{(w_t+tH)/m}$ has a uniform H\"older bound.

The extension of \cref{hessian:uniform-elliptic-extension} is uniformly elliptic and concave on all real
Hessians, with coefficient oscillation bounded by
$C|x-y|(1+\|R\|)$. The analytic estimate recalled in \cref{sec:hessian-complex-limits} therefore gives uniform $C^{2,\beta}$
bounds on a finite cover. The original smooth Hessian operator is
uniformly elliptic near the attained compact cone set. Differentiating
and applying Schauder estimates gives bounds for every higher derivative.

The linearization of
$\log(H_m(w)/\omega^n)-w-tH$ is $L_w-1$, where
\[
 L_w\xi=\frac{m\ddc\xi\wedge(\omega+\ddc w)^{m-1}
                                  \wedge\omega^{n-m}}{H_m(w)}.
\]
Its principal coefficients are positive definite. The maximum principle gives zero kernel;
smooth elliptic Fredholm theory gives surjectivity and Schauder inversion.
The implicit-function theorem proves openness. The uniform estimates
prove closedness, so the path reaches $t=1$.
For general $a$, use the metric $a\omega$ and potential $\widehat w=aw$.
Both sides of its coefficient-one equation acquire the same factor
$a^n$, so it is exactly the stated equation after division by that factor.
\end{proof}

An ordinary smooth approximation need not be admissible. To obtain an
admissible one, choose a smooth obstacle above the continuous potential
and penalize the solution's distance from that obstacle.

\begin{lemma}\label{hessian:continuous-smooth-approximation}
Every continuous admissible $\varphi$ on $X$ is the uniform limit of
smooth strictly admissible functions.
\end{lemma}
\begin{proof}
Fix a smooth $h\geq\varphi$ and a constant
\[
 b>\max\left\{1,\left\|H_m(h)/\omega^n\right\|_\infty\right\}.
\]
For each positive integer $N$ solve, by the preceding lemma,
\[
 H_m(w_N)=b e^{N(w_N-h)}\omega^n.
\]
At a positive maximum of $w_N-h$, the form $\omega+\ddc h$ lies above
the strictly admissible form $\omega+\ddc w_N$, hence is admissible there.
Monotonicity would give
$b e^{N(w_N-h)}\leq H_m(h)/\omega^n<b$, a contradiction.
Thus $w_N\leq h$. The preceding zero-order estimate also gives
\[
 \|w_N\|_\infty\leq M_N:=\|h\|_\infty+\frac{\log b}{N}.
\]

Let $x_N$ maximize $\varphi-(1+N^{-1})w_N$.
The smooth positive form
$T_N=(\omega+\ddc w_N)^{m-1}\wedge\omega^{n-m}$
defines a strictly elliptic linear operator. Admissibility gives
$(\omega+\ddc\varphi)\wedge T_N\geq0$ in distributions.
Its linear maximum principle, applied to the upper test
$(1+N^{-1})w_N$ plus the maximum value, gives at $x_N$
\[
 0\leq\bigl(\omega+(1+N^{-1})\ddc w_N\bigr)\wedge T_N.
\]
This application is to a continuous distributional subsolution of a
smooth strictly elliptic operator. It does not assume a second derivative
of $\varphi$. Writing $d_N=H_m(w_N)/\omega^n$, polarized G\aa rding
positivity yields
\[
 (N+1)d_N\geq \omega\wedge T_N/\omega^n
                         \geq d_N^{(m-1)/m}.
\]
Therefore, at this maximum point,
\[
 w_N-h\geq-E_N,\qquad
 E_N=\frac{m\log(N+1)+\log b}{N}.
\]
Since $\varphi\leq h$, the maximum of
$\varphi-(1+N^{-1})w_N$ is at most $E_N+M_N/N$.
It follows on all of $X$ that
\[
 \varphi-E_N-2M_N/N\leq w_N\leq h.
\]
Given $\delta>0$, choose an ordinary smooth majorant
$\varphi\leq h\leq\varphi+\delta$, and then $N$ so large that
$E_N+2M_N/N\leq\delta$. The resulting $w_N$ is strictly admissible
and within $\delta$ of $\varphi$. Let $\delta\downarrow0$.
No monotonicity of the chosen approximants is needed.
\end{proof}
